\chapter{Structural and Temporary Works Risk Assessment}
\label{chap:Structural and Temporary Works Risk Assessment}

%---------------------------- SECTION ----------------------------
\section{Why this assessment matters in construction}

The collapse of a temporary works structure --- formwork, falsework, temporary propping, or an improperly supported existing structure --- during the construction phase is among the most catastrophic failure modes in building and civil engineering. Unlike plant incidents or fall events, which typically produce individual injuries, temporary works collapses produce mass casualty events: the simultaneous burial or crushing of multiple operatives working within or beneath the failed structure. The BSI's Code of Practice for Temporary Works, BS 5975:2019, identifies that the majority of formwork and falsework collapses occur during concrete placing, when the imposed load is at its maximum and minor deficiencies in the temporary works design or installation that were not apparent during erection are exposed to their full design load for the first time.

The regulatory framework for temporary works in construction is clear in its structure but inconsistently applied in practice. BS 5975:2019 establishes the roles of the Temporary Works Coordinator (TWC), the Temporary Works Designer, and the Temporary Works Supervisor, and imposes a category-based classification system (T1 through T4) that determines the level of design, checking, and approval required for each temporary works scheme. Despite the clarity of this framework, many construction projects --- particularly those involving smaller contractors working on alteration and refurbishment work --- proceed without a designated TWC, without a formal temporary works design, and without a formal permit-to-load or permit-to-strike process.

The interaction between temporary works and existing structures introduces a particular complexity that is often underestimated at the assessment stage. When a contractor cuts a load-bearing wall, removes a structural beam, or excavates adjacent to an existing foundation, the load paths within the existing structure change, often in ways that are not visible and not immediately apparent. The structural survey of the existing building, the load analysis of the modified structure during the alteration phase, and the specification of temporary propping to maintain structural integrity must all be completed before work commences. Proceeding on the assumption that ``the structure has stood for years and must therefore be capable of carrying the temporary load'' is a dangerous and sometimes fatal error.

The use of proprietary temporary works systems --- Acrow props, push-pull props, system scaffolding used as falsework, adjustable steel props, and proprietary beam systems --- requires that the systems are used strictly in accordance with the manufacturer's specifications and that the imposed loads do not exceed the safe working load of the equipment. Failure to check manufacturer's guidance, and in particular failure to recognise that safe working loads reduce significantly with increasing prop extension, is a common cause of temporary works failure.

\barquote{``Temporary works do not fail because of unknown forces; they fail because known forces were not calculated, known guidance was not followed, and known responsibilities were not assigned.''}

\example{Site Reality}{A main contractor is carrying out a ground-floor extension to a Victorian terraced house. The existing party wall is load-bearing, carrying floor joists from the first floor. The contractor commences demolition of a section of the party wall to create the new opening without installing temporary propping. The labourer removing the party wall brickwork notices the floor joists above beginning to deflect as the wall is reduced below bearing level. The contractor, realising the structural implications too late, attempts to install Acrow props in an emergency. One prop is positioned off-centre beneath a rotten joist; it punches through the joist as the remaining wall section is removed, causing partial collapse of the first floor. The investigation finds that no structural engineer was appointed, no temporary works design was produced, and no consideration was given to the load-bearing function of the party wall before demolition commenced.}

%---------------------------- SECTION ----------------------------
\section{Legal and professional framework}

The primary legislative framework for structural and temporary works safety is established by the Health and Safety at Work etc.\ Act 1974, the Management of Health and Safety at Work Regulations 1999, and CDM 2015. CDM 2015 requires designers to eliminate foreseeable risks during construction through design, and to provide information about residual risks in the pre-construction information; this includes information about the structural behaviour of existing buildings during alteration and the temporary works implications of the proposed construction methodology. The Construction Phase Plan must address the arrangements for managing temporary works, including the designation of a Temporary Works Coordinator.

BS 5975:2019 (Code of Practice for Temporary Works Procedures and the Permissible Stress Design of Falsework) is the recognised British Standard for temporary works management in construction. Compliance with BS 5975 is treated by the HSE and by most principal contractors as the minimum acceptable standard for temporary works management. The Building Regulations 2010 (Part A --- Structure) impose statutory requirements on the structural adequacy of permanent works, which must not be compromised by the temporary works methodology.

\paragraph{Key frameworks relevant to this chapter include: CDM 2015; BS 5975:2019; MHSWR 1999; Building Regulations 2010 Part A; PUWER 1998.}

\begin{itemize}
  \bitem{BS 5975:2019 Code of Practice for Temporary Works}{Establishes the temporary works category system (T1--T4); defines the roles and responsibilities of the Temporary Works Coordinator, Temporary Works Designer, and Temporary Works Supervisor; requires a formal permit-to-load and permit-to-strike process for falsework; establishes the design check requirements for each category of temporary works; provides guidance on erection, use, monitoring, and dismantling.}
  \bitem{CDM 2015}{Requires designers to provide information about the structural behaviour of existing structures during the design phase; requires the Construction Phase Plan to address temporary works management arrangements; requires the principal contractor to manage and coordinate temporary works activities on site; requires contractors to assess risks from structural instability before commencement of alteration or demolition work.}
  \bitem{MHSWR 1999}{Requires suitable and sufficient risk assessment of all temporary works activities, including the assessment of structural loads during the construction phase, the adequacy of temporary supports, and the management of the permit-to-load and permit-to-strike processes.}
  \bitem{PUWER 1998}{Applies to all temporary works equipment including Acrow props, proprietary falsework systems, push-pull props, and adjustable steel beams; requires that equipment is suitable for the imposed load, used in accordance with manufacturer's instructions, and inspected before use and at appropriate intervals.}
\end{itemize}

%---------------------------- SECTION ----------------------------
\section{Conducting the assessment using the five-step method}

\subsection{Step 1 --- Identify hazards}
\paragraph{Structural and temporary works hazards include: collapse of falsework or formwork under wet concrete load during casting, causing burial of operatives beneath the fallen structure and fresh concrete; collapse of temporary propping of existing structures during alteration or demolition work, causing progressive or total structural failure; failure of temporary retaining structures during deep excavation, causing burial; failure of proprietary shoring or propping systems used at loads beyond their rated capacity or extended beyond their safe working range; premature striking of falsework before concrete has achieved sufficient strength, causing structural failure during loading; inadequate structural survey of existing buildings before alteration, causing unanticipated load path failure during the works; and failure of temporary edge protection and working platforms associated with elevated construction work.}

\subsection{Step 2 --- Decide who or what is affected}
\paragraph{Persons at risk from structural and temporary works failures include: operatives working on or beneath falsework and formwork during concrete placing, the period of highest risk; operatives involved in stripping of formwork and dismantling of falsework; tradespeople working in the vicinity of temporarily propped structures; site visitors and the site management team; members of the public adjacent to sites where temporary retaining structures are installed to support the public highway or adjacent structures; and occupants of adjacent buildings affected by ground movement or structural movement associated with temporary works operations.}

\subsection{Step 3 --- Evaluate likelihood and consequence}
\paragraph{Temporary works collapse events are typically low-frequency but very high-consequence: when they occur, they produce multiple fatalities or life-changing injuries. The inherent risk of undesigned temporary works must be rated as Very High (25/25) because the consequence of failure is catastrophic and the structural engineer's design provides the only reliable basis for safe loading. Properly designed, checked, and approved temporary works, installed in accordance with the design and monitored during loading, can reduce residual risk to Low (4--6/25). The risk assessment must differentiate between the inherent risk (no design, no check) and the residual risk (full BS 5975 compliance).}

\subsection{Step 4 --- Select and implement controls}
\paragraph{The primary control for temporary works risk is professional engineering design: appoint a structural engineer with temporary works competency to design all T2, T3, and T4 temporary works; require a design check by an independent competent person for T3 and T4 schemes; designate a Temporary Works Coordinator for all projects involving significant temporary works; implement the permit-to-load and permit-to-strike processes from BS 5975; conduct pre-use inspection of all temporary works equipment; require manufacturer's specifications for all proprietary systems and confirm that imposed loads are within rated capacity; commission a structural survey of existing buildings before any alteration work commences.}

\subsection{Step 5 --- Record and review}
\paragraph{Temporary works records must be maintained throughout the project and must include: the temporary works design calculations and drawings; the design check certificate; the permit-to-load and permit-to-strike records; the pre-use inspection records for all temporary works equipment; concrete strength test results confirming adequacy before falsework is struck. These records must be retained as part of the project health and safety file. Any event that may have affected the integrity of the temporary works --- unexpected load, impact, differential settlement, adverse weather, or reported cracking --- must trigger an immediate engineering review before work continues.}

\example{Five-Step Application}{A structural contractor is constructing a post-tensioned concrete transfer beam at first-floor level in a multi-storey residential building. The transfer beam is 900\,mm deep, 5 metres long, and carries loads from four floors above. Step 1 identifies: collapse of falsework under wet concrete load; premature striking of falsework before post-tensioning is complete; failure of proprietary beam formwork under concrete hydrostatic pressure. Step 2 identifies: four concretors on the first-floor slab; groundworkers at ground level beneath the beam; structural engineer attending for post-tensioning. Step 3 rates inherent falsework collapse risk as Very High (25/25). Step 4 implements: BS 5975 T3 temporary works design by structural engineer; independent design check; permit-to-load issued by TWC before concrete placing commences; concrete strength testing confirming 80\% design strength before post-tensioning; permit-to-strike issued by structural engineer after post-tensioning complete; PUWER inspection of all falsework components before loading. Residual risk rated Low (6/25). Step 5: daily monitoring of falsework settlement during curing.}

%---------------------------- SECTION ----------------------------
\section{Applying the hierarchy of controls}

\paragraph{Temporary works risk reduction is primarily achieved through professional design and systematic process: the hierarchy of controls at this scale is dominated by engineering design and administrative verification rather than physical guarding.}

\begin{itemize}
  \bitem{Elimination}{Design the permanent works to minimise temporary works complexity: specify precast concrete elements that do not require formwork; use proprietary permanent formwork systems that do not need to be struck; design alteration sequences that maintain load paths without the need for temporary propping.}
  \bitem{Substitution}{Replace site-assembled falsework with proprietary engineered systems with known load ratings; use precast concrete floor slabs instead of in-situ concrete where the structural design permits, to eliminate concrete-placing falsework requirements.}
  \bitem{Engineering controls}{Professional structural engineering design for all T2 and above temporary works; independent design check for T3 and above; proprietary temporary works systems used within their manufacturer's rated capacity; physical exclusion zones beneath falsework during concrete placing.}
  \bitem{Administrative controls}{Designate a Temporary Works Coordinator for all projects involving significant temporary works; implement permit-to-load and permit-to-strike processes; conduct pre-use inspection of all temporary works equipment; obtain concrete strength results before striking; commission structural survey before alteration work; provide temporary works awareness training to site supervisors.}
  \bitem{PPE}{Hard hats for all persons in the vicinity of temporary works operations; hearing protection where concrete vibration equipment is in use; safety boots and high-visibility vests; eye protection during concrete pouring.}
\end{itemize}

%---------------------------- SECTION ----------------------------
\section{Cadence of assessment and review triggers}

\paragraph{Temporary works assessments are project-phase specific and must be reviewed as each temporary works scheme is designed, installed, loaded, and struck. The assessment lifecycle mirrors the temporary works lifecycle.}

\begin{itemize}
  \bitem{Before commencement of each temporary works scheme}{A specific risk assessment for each temporary works scheme must be prepared and approved before the scheme is installed; for T2 and above, this must include the engineering design and design check.}
  \bitem{Pre-use inspection before loading}{All temporary works structures must be inspected by the Temporary Works Supervisor before the permit-to-load is issued; any discrepancy between the installed condition and the design drawing must be resolved before loading commences.}
  \bitem{During loading}{For falsework supporting concrete placing, the structure must be monitored for movement or deflection during the pour; any unplanned movement must trigger an immediate suspension of the pour and an engineering review.}
  \bitem{Before striking}{Falsework must not be struck until concrete test results confirm the required strength has been achieved and the permit-to-strike has been issued; premature striking is among the most common causes of post-pour structural failure.}
  \bitem{After any unusual event}{Any impact, unexpected settlement, flooding, frost, or significant change in the loading pattern of a temporary works structure must trigger an immediate engineering review before loading or use continues.}
\end{itemize}

%---------------------------- SECTION ----------------------------
\section{Common failure modes}

\begin{itemize}
  \bitem{No temporary works design produced}{The most fundamental failure: works proceed without a competent engineering design for the temporary structure, relying instead on the experience of the site supervisor or the site agent's judgement. Experience and judgement are not substitutes for structural calculation.}
  \bitem{No Temporary Works Coordinator designated}{Without a designated TWC, the permit-to-load and permit-to-strike processes do not occur and the oversight function required by BS 5975 is absent. The project proceeds without formal temporary works governance.}
  \bitem{Props extended beyond safe working range}{Acrow props and system props have load ratings that reduce significantly with extension. A prop correctly rated at 40\,kN at minimum extension may be rated at only 15\,kN at maximum extension. Using props at lengths beyond their rated extension is a common and potentially fatal error.}
  \bitem{Falsework struck before concrete achieves design strength}{Programme pressure drives early striking of formwork before concrete has achieved the 80\% or 100\% design strength required by the structural engineer. Concrete strength can be substantially below the characteristic strength at early ages, particularly in cold weather when curing is retarded.}
  \bitem{Structural survey not commissioned before alteration}{Proceeding with alteration of an existing building without a professional structural survey to identify load-bearing elements, hidden steelwork, and non-standard construction leads to the inadvertent removal of structural elements that were not apparent from visual inspection.}
  \bitem{Unexpected loads applied to temporary works}{Plant, materials, and operatives are positioned on or adjacent to temporary works structures without reference to the design load; the imposed load exceeds the design capacity of the temporary structure.}
\end{itemize}

\example{Failure Pattern}{A refurbishment contractor is removing a ground-floor internal wall in a three-storey office building to create an open-plan space. The wall is identified in the RAMS as ``possibly load-bearing''. No structural survey is commissioned. The contractor installs two Acrow props either side of the opening as a precaution, without a design or calculation, at an extension that places them beyond their safe working range for the estimated wall load. The RAMS is submitted and accepted by the client without scrutiny. During wall removal, the floor above begins to deflect; one of the over-extended props buckles. Emergency propping by the structural engineer appointed after the event reveals that the wall was carrying point loads from the roof structure via a series of concealed steel columns embedded in the wall thickness.}

%---------------------------- CHAPTER SUMMARY ----------------------------
\section{Chapter Summary}

\begin{table}[H]
\centering
\begin{tabularx}{\textwidth}{>{\raggedright\arraybackslash}p{3.2cm}X}
\toprule
\textbf{Assessment Element} & \textbf{Operational Requirement} \\
\midrule
Legal/Professional Basis & CDM 2015; BS 5975:2019; MHSWR 1999; Building Regulations 2010 Part A. Professional engineering design is the primary control for all T2 and above temporary works schemes. \\
Method & Apply the full five-step process and document inherent/residual risk. \\
Controls & Engineering design by competent structural engineer; designated Temporary Works Coordinator; permit-to-load and permit-to-strike processes; pre-use inspection of equipment; concrete strength confirmation before striking. \\
Cadence & Before each scheme is installed; pre-use inspection before loading; during loading for major pours; before striking; after any unusual event. \\
Assurance & Use audit, incident learning, and governance reporting to verify effectiveness. \\
\bottomrule
\end{tabularx}
\end{table}

\begin{itemize}
  \bitem{Key takeaway 1}{All T2 and above temporary works schemes require a professional structural engineering design; site supervisor experience and established practice are not substitutes for structural calculation.}
  \bitem{Key takeaway 2}{A Temporary Works Coordinator must be designated for all projects involving significant temporary works; without a TWC, the permit-to-load and permit-to-strike processes do not function.}
  \bitem{Key takeaway 3}{Falsework must not be struck until concrete strength test results confirm the required strength; programme pressure to strike early is a common cause of post-pour structural failure.}
  \bitem{Key takeaway 4}{A structural survey of existing buildings must be commissioned before any alteration work to identify load-bearing elements, hidden steelwork, and non-standard construction that may not be apparent from visual inspection.}
\end{itemize}

\barquote{In Chapter~\ref{chap:COSHH Assessment for Construction Chemicals}, we turn from this assessment to the next domain in the Prudentia framework.}
